\section{Provisioning the Network: Radio, Resources, Transport and Architecture}\label{sec:provision}

This section asks how access networks, schedulers, transport protocols and system architectures supply
the rate, delay and reliability that XR needs. Three cells hold the work: resource allocation
(\nxrNProbResourceAllocation{} papers, \nxrShProbResourceAllocation{} of the corpus), architecture and
vision (\nxrNProbArchitectureVision{}) and wireless link design (\nxrNProbWirelessLinkDesign{}).
Resource allocation is the second-largest problem cell and the fastest-growing large one: its share rose
from \nxrShProbResourceAllocationEarly{} in 2017--19 to \nxrShProbResourceAllocationLate{} in 2022--24.
Over the same periods the share of papers built on 6G rose from \nxrShTechSixGBeyondFiveGEarly{} to
\nxrShTechSixGBeyondFiveGLate{}. The three cells differ sharply in how they are evaluated.
\nxrSimResourceAllocation{} of allocation papers are simulations and \nxrRealResourceAllocation{} use a
real testbed; \nxrNoEvalArchitectureVision{} of architecture papers have no evaluation at all. Link
design is the most empirical of the three (\nxrRealWirelessLinkDesign{} real testbeds), but no paper in
any of the three cells asked users what they perceived. The argument of this section is that provisioning
research has solved well-posed optimisation problems whose inputs (QoE functions, rendering costs,
traffic and requirement figures) were assumed rather than measured. Offloading decisions are treated in
\cref{sec:compute}.

\subsection{Allocating radio, compute and money}
The design question is how to divide spectrum, power, edge GPU time and caches among XR users so that
each meets a delay or quality target. The established result is that these resources substitute for each
other. Caching rendered fields of view and computing on the device can stand in for bandwidth
\cite{W2887223659,W2945507395}, and joint uplink, downlink and edge allocation saves mobile energy in
collaborative AR compared with offloading each task on its own \cite{W2608841173}. Shared content
adds a second lever. Multicasting overlapping 360-degree viewports over mmWave raised HD delivery by
40--50\% and cut 99th-percentile delay by 17--37\% against unicast, in simulation driven by 50 real head
traces \cite{W2901851646}; the gain shrinks as viewports diverge, so it is unproven for interactive
6DoF content.

What is contested is the objective. The cell optimises at least four incompatible QoE constructs: a
utility of tracking, processing and transmission delay \cite{W2596760110}, regressions of SSIM and VMAF
\cite{W3214250514}, an attention-weighted ``Meta-Immersion'' score \cite{W4379382423} and weighted sums of
quality, delay and variation \cite{W4312555471}. Reported gains of 20--37\% are measured in each paper's
own metric and cannot be compared. Of the \nxrSynXRAQoEN{} allocation papers that optimise QoE, none was
evaluated with users. The second dispute is method. Reinforcement-learning schedulers beat greedy,
random and proportional-fair baselines \cite{W2596760110,W4381946597}, while optimisation papers reach
near-optimal decisions four orders of magnitude faster than a commercial solver \cite{W4324125331}; the
two camps rarely meet on the same traces.

Much of the rest is assumed. Rendering capacity is a free variable in joint radio--rendering allocation
\cite{W4379382423}; the only measured GPU figure is economic, a 4~GB vGPU slice that halved cloud-XR
hourly cost while holding MOS at 4.38 or above \cite{W7115818680}. Metaverse resource markets are the
extreme case. A learned double Dutch auction reached social welfare within about 5\% of the optimum
\cite{W3214250514}, and function sharing raised slice admission revenue by up to 80\% \cite{W4381946597},
but every auction, contract and incentive scheme in the cell uses synthetic valuations and arrivals.
Standards-aligned work is more realistic about traffic. It uses 3GPP XR models and PDU-set quality of
service \cite{W4410582180}, and a delay-tracking PDU-set scheduler served 4--5 times more satisfied users
than PDU-set-aware baselines under scarce resources \cite{W7168254691}. Yet only \nxrSynXRARLStdN{} of
the \nxrSynXRARLN{} reinforcement-learning allocators use these models. The few testbed studies temper
the simulated claims. On an O-RAN testbed, latency-estimation control met 10--16.67~ms targets but
over-provisioned, while packet-size prediction saved resources and lost latency control
\cite{W4410491869}; a latency-steering xApp freed 15.8\% more resources than an equivalent static slice
\cite{W4392175446}.

\subsection{The last hop: directional links, Wi-Fi and 5G}
The link question is how to keep a multi-gigabit, low-delay channel to a moving head. The cell agrees on
the diagnosis: the tail, not the peak rate, limits XR over directional links. Analysis of THz links gives
99.999\% reliability at 20~ms with guaranteed line of sight, but with blockage reliability falls to
68\% at 10~ms and 96\% at 20~ms, and tail delays reach 30--90~ms while the mean stays below 20~ms
\cite{W3022996137}. At 60~GHz, 3D head motion defeated beam search designed for 2D \cite{W2791171246},
and commercial hardware lost up to 4.72\% of packets under 45$^\circ$/s rotation \cite{W4401507876}.
Spatial diversity is the remedy that holds across studies. Dual-beam reception with coordinated access
points raised reliability by up to 22.8\% on real headset motion \cite{W4409076649}, distributed antennas
kept 12~Gb/s under full rotation \cite{W4413277491}, and steering beams from the headset's pose helps
\cite{W4313170175}: predictive beams from a neural radio field halved beam-search error on a 28~GHz array
\cite{W7204893690}.

On Wi-Fi and 5G the bound is medium access, not physical rate. With frames aligned to the 802.11ad beacon
interval, eight headsets at 400~Mbit/s met a 1~ms bound \cite{W3081368317}. Wi-Fi~7 multi-link operation
served 7 rather than 6 AR stations at 80~MHz, and an uninformed greedy policy matched congestion-aware
ones \cite{W4376480829}. Whether downlink or uplink binds is contested, and depends on the assumed rates
\cite{W4376480829,W4409221564}; on real hardware, uplink OFDMA with three headsets disrupted tracking
\cite{W4396853623}. 5G-Advanced link adaptation reports up to 67\% more satisfied XR users
\cite{W4315777456}, entirely from system-level simulation. Two technologies rest on assumption alone.
RIS-assisted XR has never run on hardware, and RIS optimisers that take seconds per slot
\cite{W4318586163} must react within a head rotation \cite{W3207401133}. No XR traffic has crossed a THz
link: the THz analyses assume Poisson arrivals and fixed-size frames \cite{W3022996137}.

\subsection{Transport over shared bottlenecks}
Transport decides what reaches the headset when the bottleneck is shared and lossy. Its share of the
corpus fell from \nxrShTechInternetTransportWebEarly{} to \nxrShTechInternetTransportWebLate{} as
delivery moved from buffered streaming to interactive rendering, although interactive traffic is less
able to absorb loss. Users confirm the risk. In cloud VR, RTT up to 90~ms was hardly noticed while loss
of up to 8\% mattered \cite{W3095344018}; with CloudXR, QoE fell at RTT of 75~ms or more or loss above
6\% (N=30) \cite{W4402156718}; with Air Link, loss raised cybersickness more than delay did (N=16)
\cite{W4404916373}. Despite this, only \nxrSynXRendFECN{} remote-rendering papers address forward error
correction or error resilience \cite{W4399340506,W4409129991}.

Congestion control for XR is nascent. L4S, designed for this traffic, appears in \nxrSynXLFourSN{}
corpus papers. An L4S pipeline for volumetric conferencing converged 45\% faster than Google congestion
control and lost no frames where WebRTC lost 3\% \cite{W4408870658}, and L4S has been carried over 5G for
split rendering \cite{W4401507827}; both studies used controlled bottlenecks with one competing scheme.
Practice has moved ahead of research. Steam Link carries video and control in one QUIC tunnel
\cite{W7160626692}, and WebRTC selective forwarding adapted 13--36\% faster than TCP for volumetric
streams in emulation \cite{W7151958829}. No paper has compared delay-based, loss-based and L4S
controllers with forward error correction on real cellular or Wi-Fi bottlenecks with cross traffic.

\subsection{Architectures and the figures they carry}
Architecture papers propose where functions sit (device, edge, cloud, ledger) and state what the network
must deliver. They agree on the qualitative list (high rate, low delay, reliability, concurrency,
interoperability) and on edge computing as a necessary tier \cite{W4312940001,W4315497913,W4391305431}.
The quantitative figures are where they fail. \Cref{tab:provision-requirements} sets stated
requirements against measured or derived figures; for rates the gap spans two to five orders of
magnitude. The cause is not error in either source but unstated operating points: an uncompressed light
field is not a compressed point cloud, and 360-degree video is not avatar state.

\begin{table}[t]
\centering\small
\caption{Stated requirements against measured or derived figures. The gaps follow from unstated content
representation, codec and control loop.}
\label{tab:provision-requirements}
\begin{tabularx}{\linewidth}{@{}p{2.3cm}XX@{}}
\toprule
Quantity & Stated & Measured or derived \\
\midrule
Per-user metaverse rate & 100--250~Mbps \cite{W4313291084,W4312940001}; $\sim$10$^3\times$ 5G
  throughput \cite{W4391305431} & Horizon Workrooms 2--3~Mbps down, 0.6~Mbps up, +0.5~Mbps per extra
  user (2--5 users) \cite{W4294643390} \\
Holographic call & 1--10~Tbps; 6.63~TBps for a light-field display \cite{W4300960574} & 5--100~Mbps
  for a compressed point cloud over live 5G \cite{W4312662428}; $<$0.2~Mbps for a driven avatar
  \cite{W7123904049} \\
Untethered VR link & 15.75--26.25~Gbps uncompressed \cite{W4401507876} & 400~Mbit/s per 4K headset,
  eight per 802.11ad AP \cite{W3081368317} \\
Cloud VR bitrate & 3GPP XR traffic model & 25\% higher effective bitrate on a Pico Neo~2
  \cite{W4399340506} \\
Interaction latency & 1~ms in \nxrSynOneMsN{} abstracts, \nxrSynOneMsUsersN{} with users & Haptic
  users tolerated 200~ms with adaptive control \cite{W4200154132}; cloud VR RTT up to 90~ms hardly
  noticed \cite{W3095344018} \\
Discomfort threshold & 30~ms \cite{W4313291084} vs $<$20~ms \cite{W4315497913} & Not derived from a
  user study in either source \\
\bottomrule
\end{tabularx}
\end{table}

The figures propagate. One survey reproduces another's requirement table with a reliability target
changed from $1-10^{-6}$ to $1-10^{-5}$ \cite{W4390204075,W4312940001}, and its research agenda reproduces
a third \cite{W4224330100}; fewer than one survey in twenty in the cell reports a search protocol. The
apparent consensus is partly one source counted several times. Measured platforms tell a different story:
throughput, latency and device load on five social VR platforms grew almost linearly with users, because
servers forward avatar data directly \cite{W4307020344,W4294643390}. The scaling bottleneck is server
logic, not access capacity, which weakens the case that 5G cannot carry the metaverse.

Two architectural claims are untested. Blockchain is placed in the interactive data path for ownership
and trust \cite{W4224330100,W4312940001}; \nxrSynXArchChainN{} architecture papers include a ledger,
\nxrSynXArchChainNoEval{} of them without any evaluation, and none has measured the delay it adds, although one concedes that replication raises delay with users
\cite{W4320717855}. Network slicing for XR, by contrast, has testbed evidence: AR, cloud gaming and video
calls were measured under RAN-slice configurations on an O-RAN 5G network \cite{W4386887692}. 6G visions
rarely map their functions onto the 5G-Advanced XR mechanisms that already exist
\cite{W4387933382,W4376607343}, so what 6G must add remains unstated.

\begin{openproblems}
\item \textbf{Allocation objectives are not validated with users.} None of the \nxrSynXRAQoEN{}
  QoE-driven allocators and no link design was evaluated with participants. A QoE-in-the-loop slicing
  study on an open 5G testbed, comparing static, latency-target and QoE-model control with ratings and
  sickness scores, would show which objective matters.
\item \textbf{Rendering costs in joint radio--GPU allocation are assumed.} Rendering capacity is a free
  variable in these models. Profiling render and encode time and GPU power on edge hardware, then
  re-running published allocators under measured costs, would show whether their rankings survive.
\item \textbf{RIS and THz have never carried XR traffic.} All RIS-for-XR work is simulated, and THz
  reliability rests on Poisson traffic. A pose-driven RIS prototype at 28 or 60~GHz, measured on frame
  deadline misses, would test the premise.
\item \textbf{Requirement figures are not derived from measured traffic.} Stated per-user rates exceed
  measured ones by up to two orders of magnitude. A traffic atlas of shared-world platforms across user
  counts and avatar fidelities, compared with 3GPP models, would replace copied tables.
\item \textbf{Transport and loss resilience are untested on real bottlenecks.} Users weigh loss above
  delay, yet L4S appears in \nxrSynXLFourSN{} papers. Comparing L4S, delay-based and loss-based control
  with forward error correction over cellular and Wi-Fi links with cross traffic, reporting p99 frame
  latency, would fill this.
\item \textbf{Blockchain latency is uncosted.} No study measures ledger delay against interaction
  deadlines. A benchmark of centralised, permissioned and rollup ledgers inside a live session would
  settle whether the ledger belongs in the data path.
\end{openproblems}
